\documentclass[11pt,openany]{report}
\usepackage[letterpaper,margin=0.82in,headheight=14pt]{geometry}
\usepackage[T1]{fontenc}
\usepackage[utf8]{inputenc}
\usepackage{lmodern}
\usepackage{microtype}
\usepackage{amsmath,amssymb,amsthm,mathtools}
\usepackage{booktabs,longtable,tabularx,array,multirow}
\usepackage{graphicx,float,placeins}
\usepackage[dvipsnames]{xcolor}
\usepackage{enumitem}
\usepackage{fancyhdr}
\usepackage{listings}
\usepackage{url}
\usepackage[hidelinks]{hyperref}
\usepackage[nameinlink,noabbrev]{cleveref}
\usepackage{needspace}
\allowdisplaybreaks

\hypersetup{
  pdftitle={Projector-Generated Special Functions III},
  pdfauthor={Komorebi P.},
  pdfsubject={Squarefree Stieltjes-Jacobi functions, Laguerre packet universality, Hankel determinants, and moment-exact reduction}
}

\definecolor{Ink}{HTML}{17202A}
\definecolor{Blue}{HTML}{2E86AB}
\definecolor{Teal}{HTML}{2A9D8F}
\definecolor{Orange}{HTML}{F4A261}
\definecolor{Magenta}{HTML}{A23E8C}
\definecolor{Red}{HTML}{E76F51}
\definecolor{Paper}{HTML}{FAFAF7}
\definecolor{Muted}{HTML}{5D6D7E}
\definecolor{CodeBG}{HTML}{F3F5F5}

\pagecolor{Paper}
\color{Ink}
\graphicspath{{figures/}}
\setlength{\parindent}{0pt}
\setlength{\parskip}{0.52em}
\raggedbottom

\pagestyle{fancy}
\fancyhf{}
\fancyhead[L]{\small Projector-Generated Special Functions III}
\fancyhead[R]{\small Komorebi P.}
\fancyfoot[C]{\thepage}
\renewcommand{\headrulewidth}{0.4pt}

\newtheorem{theorem}{Theorem}[chapter]
\newtheorem{proposition}[theorem]{Proposition}
\newtheorem{lemma}[theorem]{Lemma}
\newtheorem{corollary}[theorem]{Corollary}
\newtheorem{conjecture}[theorem]{Conjecture}
\theoremstyle{definition}
\newtheorem{definition}[theorem]{Definition}
\newtheorem{example}[theorem]{Worked example}
\newtheorem{algorithm}[theorem]{Algorithm}
\theoremstyle{remark}
\newtheorem{remark}[theorem]{Remark}
\newtheorem{boundary}[theorem]{Boundary statement}

\newcommand{\C}{\mathbb C}
\newcommand{\R}{\mathbb R}
\newcommand{\N}{\mathbb N}
\newcommand{\Z}{\mathbb Z}
\newcommand{\e}{\mathrm e}
\newcommand{\ii}{\mathrm i}
\newcommand{\Ecal}{\mathcal E}
\newcommand{\Jcal}{\mathcal J}
\newcommand{\Kcal}{\mathcal K}
\newcommand{\Mcal}{\mathcal M}
\newcommand{\Pcal}{\mathcal P}
\newcommand{\norm}[1]{\left\lVert#1\right\rVert}
\newcommand{\abs}[1]{\left\lvert#1\right\rvert}
\newcommand{\dd}{\,\mathrm d}
\newcommand{\diag}{\operatorname{diag}}
\newcommand{\supp}{\operatorname{supp}}
\newcommand{\Res}{\operatorname{Res}}
\newcommand{\Var}{\operatorname{Var}}
\newcommand{\Prob}{\mathbb P}
\newcommand{\Expect}{\mathbb E}
\newcommand{\ket}[1]{\left|#1\right\rangle}
\newcommand{\bra}[1]{\left\langle#1\right|}
\newcommand{\braket}[2]{\left\langle#1\middle|#2\right\rangle}
\newcommand{\inner}[2]{\left\langle#1,#2\right\rangle}

\newcolumntype{L}[1]{>{\raggedright\arraybackslash}p{#1}}

\lstdefinestyle{pythonstyle}{
  language=Python,
  basicstyle=\ttfamily\scriptsize,
  backgroundcolor=\color{CodeBG},
  frame=single,
  rulecolor=\color{Muted!45},
  breaklines=true,
  breakatwhitespace=false,
  showstringspaces=false,
  numbers=left,
  numberstyle=\tiny\color{Muted},
  numbersep=7pt,
  tabsize=4,
  columns=fullflexible,
  keepspaces=true
}
\lstset{style=pythonstyle}

\newcommand{\FigureWide}[3]{%
\Needspace{0.53\textheight}
\begin{figure}[H]
\centering
\includegraphics[width=0.92\textwidth,height=0.44\textheight,keepaspectratio]{#1.pdf}
\caption{#2 #3}
\label{fig:#1}
\end{figure}
}

\newcommand{\InlinePart}[1]{%
\FloatBarrier
\Needspace{10\baselineskip}
\refstepcounter{part}%
\addcontentsline{toc}{part}{\protect\numberline{\thepart}#1}%
\vspace{1.4em}
{\color{Blue}\Large\bfseries Part \thepart\par}
\vspace{0.2em}
{\color{Blue}\LARGE\bfseries #1\par}
\vspace{0.9em}
}

\newcommand{\FlowChapter}[1]{%
\FloatBarrier
\Needspace{14\baselineskip}
\refstepcounter{chapter}%
\setcounter{section}{0}%
\setcounter{subsection}{0}%
\addcontentsline{toc}{chapter}{\protect\numberline{\thechapter}#1}%
\markboth{#1}{#1}%
\vspace{1.3em}
{\color{Muted}\large\bfseries Chapter \thechapter\par}
\vspace{0.15em}
{\LARGE\bfseries #1\par}
\vspace{0.9em}
}

\title{\textbf{Projector-Generated Special Functions III}\\[0.25em]
\Large Squarefree Stieltjes--Jacobi Functions, Laguerre Packet Universality, Hankel Determinants, and Moment-Exact Reduction}
\author{Komorebi P.}
\date{}

\begin{document}
\maketitle

\begin{abstract}
This paper develops a third family of special functions generated by the projector-resolved calculus. The starting point is the positive squarefree packet energy introduced in the preceding paper. Its values are reinterpreted as the moments of a discrete measure on the unit interval. That measure produces a Stieltjes resolvent, an infinite Jacobi operator, a continued fraction, positive rational approximants, Hankel determinants, Christoffel--Darboux kernels, and Gaussian quadrature rules. These are not independent definitions. They are exact coordinate realizations of one positive packet seed.

The moment sequence is proved to be strictly totally positive in the generalized Hankel sense. Its resolvent is analytic and zero-free on the slit plane $\C\setminus[1,\infty)$. Lanczos reduction produces a finite Jacobi model that preserves the first $2m$ moments exactly after $m$ steps. Every finite model has positive weights, real nodes in the packet spectrum, and a rational resolvent with poles confined to the physical cut.

A second layer appears after exponential tilting by the packet depth $k$. If
\[
\alpha=\frac{2a-1}{\beta},
\]
then the scaled endpoint coordinate $u=k(1-x)$ converges to the Gamma law with shape $\alpha$. Fixed-degree orthogonal packet polynomials converge to monic generalized Laguerre polynomials. The Jacobi coefficients, Christoffel--Darboux kernel, Gaussian quadrature nodes, and Gaussian quadrature weights converge to their Gamma--Laguerre counterparts. The Gamma--Laguerre operator is therefore the universal local fixed point of the squarefree packet family.

The paper also proves a complete leading asymptotic for shifted Hankel determinants. For the principal parameters $a=\beta=2$, the energy resolvent has a square-root cusp at the endpoint of its cut with an explicit constant. The positive theory is kept separate from the unresolved signed M\"obius coherence problem. The latter remains the arithmetic obstruction carrying the Riemann hypothesis.

All figures were produced by the accompanying Python program. The release contains 22 vector figures, matching PNG files, exact tables, 34 built-in checks, 19 independent tests, and the complete Python source in the accompanying release.
\end{abstract}

\chapter*{Status and claim boundary}
\addcontentsline{toc}{chapter}{Status and claim boundary}
The paper proves the positive moment representation, strict generalized Hankel positivity, the zero-free Stieltjes resolvent theorem, the Jacobi continued fraction, moment-exact finite reduction, Gamma scaling, fixed-degree Laguerre convergence, recurrence-coefficient convergence, Christoffel--Darboux convergence, Gaussian quadrature convergence, and the shifted Hankel determinant asymptotic under the hypotheses stated below.

Classical inputs include squarefree counting, partial summation, the compact moment problem, orthogonal polynomial recurrence theory, Gaussian quadrature, Lanczos tridiagonalization, and generalized Laguerre polynomials. The construction of the packet measure and its arithmetic energy law comes from the second paper of this series. The logarithmic number-basis Hamiltonian discussed near the end is compatible with the engineered quantum setting developed by Wei and collaborators~\cite{Wei2026}.

The positive packet theory does not prove RH. It controls the diagonal squarefree energy and supplies canonical finite models. RH remains equivalent to a subpolynomial bound for a signed M\"obius coherence ratio. No positivity statement proved here implies that missing signed estimate.

\tableofcontents
\listoffigures

\InlinePart{The squarefree packet measure}

\FlowChapter{One positive seed and its transforms}

Fix parameters
\[
a>\frac12,
\qquad
\beta>0,
\qquad
q\in\N.
\]
Let $\chi$ be a Dirichlet character of conductor $q$. Since $\abs{\chi(n)}^2$ is the indicator of $(n,q)=1$, define
\[
w_n
=
\frac{\mu(n)^2\abs{\chi(n)}^2}{n^{2a}}
\]
and
\[
x_n
=
\left(1-\frac1{n^\beta}\right)^2.
\]
Every $w_n$ is nonnegative and every node belongs to $[0,1)$. The packet measure is
\[
\boxed{
\nu_{a,\beta,q}
=
\sum_{n\ge1}w_n\,\delta_{x_n}.
}
\]
Its moments are
\[
\boxed{
E_{a,\beta,q}(k)
=
\int_0^1x^k\dd\nu(x)
=
\sum_{n\ge1}
\frac{\mu(n)^2\abs{\chi(n)}^2}{n^{2a}}
\left(1-\frac1{n^\beta}\right)^{2k}.
}
\]
The condition $a>1/2$ makes the total mass finite.

\FigureWide{01_theory_pipeline}{The theory generated by one positive packet seed.}{The measure gives moments, a resolvent, a Jacobi operator, finite quadrature models, tilted endpoint laws, and a Laguerre fixed point.}

\begin{definition}[Squarefree packet moment function]
The map
\[
(k,a,\beta,q)
\longmapsto
E_{a,\beta,q}(k)
\]
is the squarefree packet moment function. For fixed parameters it is a sequence in $k$. For fixed $k$ it is a positive special value built from the squarefree Dirichlet series.
\end{definition}

\begin{proposition}[Exact finite-difference zeta-ratio coordinate]
Let
\[
D_q(s)
=
\sum_{\substack{n\ge1\\(n,q)=1}}
\frac{\mu(n)^2}{n^s}.
\]
Then
\[
D_q(s)
=
\frac{\zeta(s)}{\zeta(2s)}
\prod_{p\mid q}(1+p^{-s})^{-1}
\]
and
\[
\boxed{
E_{a,\beta,q}(k)
=
\sum_{j=0}^{2k}
(-1)^j\binom{2k}{j}
D_q(2a+\beta j).
}
\]
\end{proposition}

\begin{proof}
The Euler factor of $\mu(n)^2$ at a prime $p\nmid q$ is $1+p^{-s}$. At a prime dividing $q$, the nontrivial term is removed. This gives the formula for $D_q$. Expanding $(1-n^{-\beta})^{2k}$ by the binomial theorem and interchanging the finite sum with the absolutely convergent Dirichlet series gives the second identity.
\end{proof}

\begin{example}[Principal values]
For $a=\beta=2$ and $q=1$,
\[
E_k
=
\sum_{j=0}^{2k}
(-1)^j\binom{2k}{j}
\frac{\zeta(4+2j)}{\zeta(8+4j)}.
\]
The Python program checks this formula against the direct squarefree sum for $0\le k\le5$.
\end{example}

\FigureWide{02_squarefree_packet_measure}{The discrete packet measure on $[0,1)$.}{Squarefree integers produce positive atoms whose locations accumulate at the spectral endpoint $x=1$. Their weights decay as $n^{-2a}$.}

\FlowChapter{Positive contraction realization}

Let
\[
\mathcal H_q
=
\ell^2\{n\ge1:(n,q)=1,\ \mu(n)^2=1\}.
\]
Define the diagonal contraction
\[
T\ket n=x_n\ket n
\]
and the seed vector
\[
\ket v
=
\sum_n\sqrt{w_n}\ket n.
\]
Then $0\le T<I$ and
\[
\boxed{
E_k
=
\bra vT^k\ket v.
}
\]
The complete positive theory begins with this one matrix element.

\begin{theorem}[Hausdorff moment structure]
The sequence $(E_k)_{k\ge0}$ is a Hausdorff moment sequence. For every $m,k\ge0$,
\[
\boxed{
(-1)^m\Delta^mE_k
=
\sum_nw_nx_n^k(1-x_n)^m
\ge0.
}
\]
\end{theorem}

\begin{proof}
The identity follows from
\[
\Delta^m x^k=x^k(x-1)^m.
\]
Multiplication by $(-1)^m$ gives $x^k(1-x)^m$. Integration against the positive measure proves the claim.
\end{proof}

\begin{corollary}
The sequence is decreasing, convex in the alternating finite-difference sense, and log-convex:
\[
E_k^2\le E_{k-1}E_{k+1}.
\]
\end{corollary}

\begin{proof}
Only log-convexity needs comment. Apply Cauchy--Schwarz to the functions $x^{(k-1)/2}$ and $x^{(k+1)/2}$ in $L^2(\nu)$.
\end{proof}

\FigureWide{03_positive_contraction}{The operator realization of the moments.}{The arithmetic weights form the seed vector. The diagonal contraction supplies the packet depth. Their matrix element is exactly the energy.}

\input{generated_tables/energy.tex}

\FigureWide{04_energy_decay}{Energy decay in the principal case.}{The computed values approach the constant $3/(2^{5/2}\pi^{3/2})$ multiplying $k^{-3/2}$.}

\InlinePart{Total positivity and the Stieltjes-Jacobi function}

\FlowChapter{Strict generalized Hankel positivity}

For increasing integer sequences
\[
0\le r_0<\cdots<r_d,
\qquad
0\le s_0<\cdots<s_d,
\]
define the generalized Hankel minor
\[
\Delta_{\mathbf r,\mathbf s}
=
\det\left[E_{r_i+s_j}\right]_{i,j=0}^{d}.
\]

\begin{theorem}[Strict total positivity]
Every generalized Hankel minor is strictly positive:
\[
\boxed{
\Delta_{\mathbf r,\mathbf s}>0.
}
\]
\end{theorem}

\begin{proof}
Use the discrete Andr\'eief identity to write the determinant as a sum over distinct support nodes:
\[
\Delta_{\mathbf r,\mathbf s}
=
\frac1{(d+1)!}
\sum_{n_0,\ldots,n_d}
\left(\prod_{j=0}^{d}w_{n_j}\right)
\det[x_{n_j}^{r_i}]
\det[x_{n_j}^{s_i}].
\]
Terms with repeated nodes vanish. Order the remaining nodes so that
\[
0<x_{n_0}<\cdots<x_{n_d}<1.
\]
A generalized Vandermonde determinant with increasing nonnegative exponents is positive on strictly increasing positive nodes. Both determinants therefore have the same positive sign. The support contains infinitely many positive nodes with positive weights, so at least one term is strictly positive.
\end{proof}

\begin{corollary}[Strict Hankel positivity]
For every $d\ge0$, the matrix
\[
H_d=(E_{i+j})_{i,j=0}^{d}
\]
is positive definite.
\end{corollary}

\begin{proof}
For $c\ne0$,
\[
c^*H_dc
=
\int_0^1
\left|\sum_{j=0}^{d}c_jx^j\right|^2\dd\nu(x)>0.
\]
A nonzero polynomial cannot vanish on the infinite support of $\nu$.
\end{proof}

\FigureWide{05_total_positivity}{Selected generalized Hankel minors.}{The heatmap records base-ten logarithms of exact high-precision minors. The proof covers every increasing pair of exponent sets.}

\FlowChapter{The zero-free Stieltjes-Jacobi function}

Define the packet-energy resolvent
\[
\boxed{
\Ecal_{a,\beta,q}(z)
=
\sum_{k=0}^{\infty}E_kz^k
=
\int_0^1\frac{\dd\nu(x)}{1-zx}.
}
\]
The power series converges for $\abs z<1$. The integral continues it to the slit plane.

\begin{theorem}[Zero-free slit-plane continuation]
The function $\Ecal_{a,\beta,q}$ is analytic and zero-free on
\[
\boxed{
\C\setminus[1,\infty).
}
\]
\end{theorem}

\begin{proof}
For real $z<1$, every integrand $(1-zx)^{-1}$ is positive, hence $\Ecal(z)>0$. If $\operatorname{Im}z>0$, then
\[
\operatorname{Im}\frac1{1-zx}
=
\frac{x\operatorname{Im}z}{\abs{1-zx}^2}
\ge0.
\]
The measure has positive mass on $x>0$, so the integral has strictly positive imaginary part. For $\operatorname{Im}z<0$, the imaginary part is strictly negative. A zero is impossible in all three regions.
\end{proof}

\begin{remark}
This is a proved zero-free function family generated by the calculus. Its zero-free property comes from a positive packet measure and not from a conjectural arithmetic cancellation.
\end{remark}

Normalize the measure by $\widehat\nu=\nu/E_0$ and let $p_m$ be its orthonormal polynomials. They obey
\[
xp_m(x)
=b_{m+1}p_{m+1}(x)+a_mp_m(x)+b_mp_{m-1}(x),
\]
where $b_m>0$. The Jacobi matrix is
\[
\Jcal=
\begin{pmatrix}
a_0&b_1&0&\cdots\\
b_1&a_1&b_2&\ddots\\
0&b_2&a_2&\ddots\\
\vdots&\ddots&\ddots&\ddots
\end{pmatrix}.
\]
The normalized Stieltjes transform
\[
S(w)=\int\frac{\dd\widehat\nu(x)}{w-x}
\]
has the Jacobi continued fraction
\[
\boxed{
S(w)
=
\cfrac{1}{w-a_0-
\cfrac{b_1^2}{w-a_1-
\cfrac{b_2^2}{w-a_2-\cdots}}}.
}
\]
Since
\[
\Ecal(z)=\frac{E_0}{z}S(1/z),
\]
this continued fraction is an exact representation of the packet-energy resolvent.

\FigureWide{06_edge_cusp}{The spectral endpoint in the principal case.}{The resolvent remains finite at $z=1$ but approaches it through a square-root cusp with a computable constant.}

\FigureWide{07_global_jacobi}{The global packet Jacobi matrix.}{Lanczos tridiagonalization turns the infinite diagonal number-basis operator into one canonical nearest-neighbor transfer system.}

\input{generated_tables/jacobi.tex}

\FigureWide{08_continued_fraction}{The Jacobi continued fraction.}{Every coefficient is determined by the positive packet measure. Finite truncations preserve positivity and produce rational reduced models.}

\InlinePart{Finite positive models}

\FlowChapter{Gauss packets and Pad\'e-type resolvents}

Let $J_m$ denote the leading $m\times m$ Jacobi matrix. Its eigenvalues are
\[
0\le x_1^{(m)}<\cdots<x_m^{(m)}<1.
\]
If $u_j$ is a normalized eigenvector, define
\[
\lambda_j^{(m)}=\abs{(u_j)_1}^2.
\]
Then every weight is positive and their sum is one.

\begin{theorem}[Moment-exact packet reduction]
For every polynomial $P$ of degree at most $2m-1$,
\[
\boxed{
\int P(x)\dd\nu(x)
=
E_0\sum_{j=1}^{m}\lambda_j^{(m)}P(x_j^{(m)}).
}
\]
Equivalently,
\[
\boxed{
E_r
=
E_0e_1^TJ_m^re_1,
\qquad
0\le r\le2m-1.
}
\]
\end{theorem}

\begin{proof}
This is Gaussian quadrature for the positive measure $\widehat\nu$. It can also be obtained directly from Lanczos. The Krylov space
\[
\operatorname{span}\{v,Tv,\ldots,T^{m-1}v\}
\]
contains every polynomial in $T$ of degree below $m$ applied to $v$. Orthogonality of the Lanczos residual makes the reduced bilinear form exact through degree $2m-1$.
\end{proof}

\begin{corollary}
The reduced resolvent
\[
\Ecal_m(z)
=
E_0\sum_{j=1}^{m}
\frac{\lambda_j^{(m)}}{1-zx_j^{(m)}}
\]
has positive residues and poles only at
\[
\frac1{x_j^{(m)}}\in(1,\infty).
\]
Its Taylor expansion agrees with $\Ecal(z)$ through degree $2m-1$.
\end{corollary}

\FigureWide{09_pade_approximants}{Positive rational approximants to the resolvent.}{The finite Jacobi models reproduce low-order moments exactly and converge rapidly on the tested compact real interval below the cut.}

\input{generated_tables/pade.tex}

\FigureWide{10_poles_and_weights}{Poles and weights of a ten-node packet model.}{All nonzero poles lie on the physical cut and every spectral weight is positive. A zero node represents the endpoint atom and produces no finite pole.}

\FigureWide{11_moment_matching}{Exact moment matching.}{An $m$-node model agrees through degree $2m-1$. The first unmatched moments expose the boundary of the finite contract.}

\FigureWide{18_moment_exact_reduction}{The reduction map.}{The finite Jacobi operator is not an arbitrary fit. It is the unique Gaussian packet model attached to the Krylov basis.}

\InlinePart{Gamma tilting and Laguerre universality}

\FlowChapter{Endpoint scaling of the packet measure}

Set
\[
\alpha=\frac{2a-1}{\beta}
\]
and define the tilted probability measure
\[
\dd\nu_k(x)
=
\frac{x^k}{E_k}\dd\nu(x).
\]
The relevant endpoint coordinate is
\[
\boxed{
u=k(1-x).
}
\]
For the atom $x_n=(1-n^{-\beta})^2$,
\[
u
=k\left(2n^{-\beta}-n^{-2\beta}\right).
\]

Let
\[
Q_q(X)
=
\sum_{\substack{n\le X\\(n,q)=1}}\mu(n)^2.
\]
The squarefree counting law is
\[
Q_q(X)=\delta_qX+O_q(\sqrt X),
\qquad
\delta_q=\frac1{\zeta(2)}\prod_{p\mid q}\frac{p}{p+1}.
\]

\begin{theorem}[Gamma packet scaling]
Under $\nu_k$, the coordinate $u=k(1-x)$ converges in distribution to the Gamma law with shape $\alpha$ and scale one. For every bounded continuous $\varphi$,
\[
\boxed{
\lim_{k\to\infty}
\int\varphi(k(1-x))\dd\nu_k(x)
=
\frac1{\Gamma(\alpha)}
\int_0^\infty
\varphi(u)u^{\alpha-1}\e^{-u}\dd u.
}
\]
\end{theorem}

\begin{proof}
Write the numerator as
\[
\sum_{\substack{n\ge1\\(n,q)=1}}
\mu(n)^2n^{-2a}(1-n^{-\beta})^{2k}
\varphi\left(k[1-(1-n^{-\beta})^2]\right).
\]
Put $n=(2k)^{1/\beta}y$. Partial summation with the squarefree counting law replaces the sum, after multiplication by $(2k)^\alpha$, by
\[
\delta_q\int_0^\infty
y^{-2a}\e^{-y^{-\beta}}
\varphi(y^{-\beta})\dd y.
\]
The substitution $u=y^{-\beta}$ gives
\[
\frac{\delta_q}{\beta}
\int_0^\infty
\varphi(u)u^{\alpha-1}\e^{-u}\dd u.
\]
Taking $\varphi=1$ gives the denominator asymptotic. Division yields the normalized Gamma law. Standard truncation at small and large $u$ supplies domination.
\end{proof}

\begin{corollary}[Energy asymptotic]
\[
\boxed{
E_{a,\beta,q}(k)
\sim
A_{a,\beta,q}k^{-\alpha},
\qquad
A_{a,\beta,q}
=
\frac{\delta_q}{\beta}
\Gamma(\alpha)2^{-\alpha}.
}
\]
\end{corollary}

\begin{example}[Principal Gamma law]
For $a=\beta=2$ and $q=1$,
\[
\alpha=\frac32,
\qquad
A=\frac{3}{2^{5/2}\pi^{3/2}}.
\]
Thus
\[
E_k\sim\frac{3}{2^{5/2}\pi^{3/2}}k^{-3/2}.
\]
The scaled variable $k(1-x_n)$ is asymptotically the same as $2k/n^2$.
\end{example}

\FigureWide{12_gamma_tilt}{Convergence of the tilted packet law.}{The discrete packet histograms approach the Gamma density. Arithmetic oscillations remain visible at finite depth and decay under scaling.}

\FlowChapter{Laguerre packet polynomials}

Let $P_m^{(k)}$ be the monic orthogonal polynomial of degree $m$ for the pushed-forward law of $u=k(1-x)$ under $\nu_k$. Define the monic Gamma--Laguerre polynomial
\[
\boxed{
\Lambda_m^{(\alpha)}(u)
=(-1)^m m!L_m^{(\alpha-1)}(u).
}
\]
Its squared norm under the normalized Gamma law is
\[
\boxed{
h_m^{(\alpha)}=m!(\alpha)_m.
}
\]

\begin{theorem}[Fixed-degree Laguerre universality]
For every fixed degree $m$,
\[
\boxed{
P_m^{(k)}(u)
\longrightarrow
\Lambda_m^{(\alpha)}(u)
}
\]
coefficientwise and locally uniformly as $k\to\infty$.
\end{theorem}

\begin{proof}
The Gamma scaling theorem gives convergence of every fixed moment of the pushed-forward measures. The Gamma moment matrix of order $m$ is positive definite. The coefficients of the monic orthogonal polynomial are the solution of a linear system whose entries are moments. Continuity of the solution under convergence of the positive definite moment matrix gives coefficientwise convergence. Local uniform convergence follows because the degree is fixed.
\end{proof}

\FigureWide{13_laguerre_polynomial_limit}{Polynomial convergence in degrees one through four.}{The packet polynomials arise from the exact discrete arithmetic measure. Their fixed-degree scaling limit is the classical monic Laguerre family.}

The monic Laguerre recurrence is
\[
\Lambda_{m+1}(u)
=
(u-2m-\alpha)\Lambda_m(u)
-m(m+\alpha-1)\Lambda_{m-1}(u).
\]
The orthonormal Jacobi coefficients are
\[
a_m^{(\infty)}=2m+\alpha,
\qquad
b_{m+1}^{(\infty)}=\sqrt{(m+1)(m+\alpha)}.
\]

\begin{corollary}[Jacobi fixed point]
For every fixed $m$,
\[
\boxed{
a_m(k)\to2m+\alpha
}
\]
and
\[
\boxed{
b_{m+1}(k)\to\sqrt{(m+1)(m+\alpha)}.
}
\]
Every fixed leading block of the scaled tilted Jacobi matrix converges to the Gamma--Laguerre Jacobi matrix.
\end{corollary}

\FigureWide{14_recurrence_coefficients}{Convergence of the packet recurrence coefficients.}{The finite arithmetic Jacobi systems approach one universal Gamma--Laguerre transfer operator determined only by $\alpha=(2a-1)/\beta$.}

\input{generated_tables/recurrence.tex}

\InlinePart{Determinants, kernels, and quadrature limits}

\FlowChapter{Shifted Hankel determinant asymptotics}

For $d\ge0$, define
\[
\Delta_d(k)
=
\det\left[E_{k+i+j}\right]_{i,j=0}^{d}.
\]
Let
\[
G_d(\alpha)
=
\prod_{j=0}^{d}j!(\alpha)_j.
\]

\begin{theorem}[Complete leading Hankel asymptotic]
For fixed $d$,
\[
\boxed{
\Delta_d(k)
\sim
A^{d+1}G_d(\alpha)
\,k^{-(d+1)\alpha-d(d+1)}
}
\]
with
\[
A=\frac{\delta_q}{\beta}\Gamma(\alpha)2^{-\alpha}.
\]
\end{theorem}

\begin{proof}
Factor $E_k$ from every row of the determinant:
\[
\Delta_d(k)
=E_k^{d+1}
\det\left[
\int x^{i+j}\dd\nu_k(x)
\right]_{i,j=0}^{d}.
\]
Apply the affine change $x=1-u/k$. The determinant of the triangular basis change from $1,x,\ldots,x^d$ to $1,u,\ldots,u^d$ contributes
\[
k^{-d(d+1)}.
\]
The remaining determinant is the Hankel determinant of the first $2d$ moments of the pushed-forward $u$ law. Those moments converge to the Gamma moments $(\alpha)_r$. The Gamma Hankel determinant equals the product of the monic orthogonal polynomial norms,
\[
\det[(\alpha)_{i+j}]_{i,j=0}^{d}
=
\prod_{j=0}^{d}j!(\alpha)_j.
\]
Finally, $E_k^{d+1}\sim A^{d+1}k^{-(d+1)\alpha}$.
\end{proof}

\FigureWide{15_hankel_determinants}{Scaled shifted Hankel determinants.}{Four determinant orders converge to the Gamma--Laguerre constants predicted by the theorem.}

\input{generated_tables/hankel.tex}

\FlowChapter{Christoffel--Darboux universality}

Let $p_m^{(k)}$ denote the orthonormal packet polynomial in the scaled coordinate $u$ and define
\[
K_d^{(k)}(u,v)
=
\sum_{m=0}^{d}p_m^{(k)}(u)p_m^{(k)}(v).
\]
The Gamma--Laguerre kernel is
\[
K_d^{(\alpha)}(u,v)
=
\sum_{m=0}^{d}
\frac{\Lambda_m^{(\alpha)}(u)\Lambda_m^{(\alpha)}(v)}{m!(\alpha)_m}.
\]

\begin{theorem}[Fixed-degree kernel convergence]
For every fixed $d$,
\[
\boxed{
K_d^{(k)}(u,v)
\longrightarrow
K_d^{(\alpha)}(u,v)
}
\]
locally uniformly on $\C^2$.
\end{theorem}

\begin{proof}
The monic polynomials and their norms converge degree by degree. The kernel is a finite sum, so local uniform convergence follows.
\end{proof}

\FigureWide{16_christoffel_darboux}{Christoffel--Darboux kernel convergence.}{The packet kernel, the Gamma--Laguerre kernel, and their difference are shown at a large but finite tilt depth.}

\FlowChapter{Gauss--Laguerre packet quadrature}

The $m$-node Gaussian rule for the scaled packet law has nodes equal to the eigenvalues of the $m\times m$ tilted Jacobi matrix and weights equal to the squares of the first eigenvector components.

\begin{theorem}[Quadrature fixed point]
For fixed $m$, the packet quadrature nodes and weights converge to the normalized Gauss--Laguerre nodes and weights for the density
\[
\frac{u^{\alpha-1}\e^{-u}}{\Gamma(\alpha)}.
\]
\end{theorem}

\begin{proof}
The finite Jacobi matrices converge entrywise. Their limiting spectra are simple. Eigenvalue and normalized spectral-weight continuity gives convergence of nodes and weights.
\end{proof}

\FigureWide{17_gauss_laguerre_limit}{Gauss quadrature convergence.}{Node and weight errors contract toward zero as the arithmetic packet tilt approaches its Gamma fixed point.}

\input{generated_tables/quadrature.tex}

\InlinePart{The spectral edge and the universal family}

\FlowChapter{A square-root cusp in the principal resolvent}

For the principal parameters $a=\beta=2$ and $q=1$, the atom locations satisfy
\[
1-x_n
=
\frac2{n^2}-\frac1{n^4}.
\]
Near the endpoint, the packet measure has density equivalent
\[
\dd\nu
\sim
C y^{1/2}\dd y,
\qquad
y=1-x,
\qquad
C=\frac{3}{2^{3/2}\pi^2}.
\]

\begin{theorem}[Principal edge cusp]
As $z\uparrow1$ through real values,
\[
\boxed{
\Ecal(1)-\Ecal(z)
\sim
\frac{3}{2^{3/2}\pi}\sqrt{1-z}.
}
\]
\end{theorem}

\begin{proof}
Put $\varepsilon=1-z$. The singular part is
\[
\int_0^\eta
C y^{1/2}
\left(
\frac1y-
\frac1{y+\varepsilon}
\right)\dd y.
\]
After $y=\varepsilon u$, the leading term is
\[
C\varepsilon^{1/2}
\int_0^\infty
\frac{u^{-1/2}}{1+u}\dd u
=
C\pi\varepsilon^{1/2}.
\]
Since $C\pi=3/(2^{3/2}\pi)$, the formula follows. Contributions away from the endpoint are $O(\varepsilon)$.
\end{proof}

\FlowChapter{Universality classes}

The local fixed point depends on the original arithmetic parameters only through
\[
\boxed{
\alpha=\frac{2a-1}{\beta}.
}
\]
Two packet systems with the same $\alpha$ have the same Gamma limit, Laguerre recurrence, determinant constants after energy normalization, fixed-degree kernel, and quadrature limit.

\FigureWide{19_parameter_universality}{The parameter map.}{Curves of constant $\alpha$ define universality classes. The principal RH energy occupies the class $\alpha=3/2$.}

\begin{example}[Three parameter pairs in one class]
The choices
\[
(a,\beta)=(2,2),
\qquad
\left(\frac54,1\right),
\qquad
\left(\frac72,4\right)
\]
all give $\alpha=3/2$. Their unscaled measures differ, but the endpoint packet polynomials converge to the same Laguerre family.
\end{example}

\FigureWide{20_dependency_map}{Logical dependency of the results.}{The positive measure supports two branches. One leads to total positivity and resolvent reduction. The other leads to Gamma scaling and Laguerre universality.}

\FigureWide{22_function_family_map}{The third generated function family.}{The resolvent, Jacobi function, rational models, Hankel determinants, Laguerre fixed point, and quadrature rules are exact transforms of one packet measure.}

\InlinePart{Quantum coordinates and the arithmetic boundary}

\FlowChapter{Number-basis realization}

On the squarefree number basis define
\[
H\ket n=\log n\ket n
\]
and define the positive contraction directly by
\[
T\ket n
=
\left(1-n^{-\beta}\right)^2\ket n.
\]
On the invariant subspace spanned by $n\ge2$, introduce
\[
K_\beta\ket n
=-\log(1-n^{-\beta})\ket n.
\]
Then
\[
T=\e^{-2K_\beta}
\]
on that subspace.  The remaining basis vector $\ket1$ is the zero-eigenvalue packet of $T$ and is kept separately, avoiding the singular logarithm at $n=1$.
For
\[
\ket{v_a}
=
\sum_{\substack{n\ge1\\(n,q)=1}}
\frac{\mu(n)^2}{n^a}\ket n,
\]
the packet moments satisfy
\[
E_{a,\beta,q}(k)
=
\bra{v_a}T^k\ket{v_a}.
\]
The operators $H$, $K_\beta$, and $T$ commute because they are functions of the number operator.

The logarithmic Hamiltonian $H\ket n=\log n\ket n$ is the same spectral coordinate used in the engineered quantum systems of Wei and collaborators~\cite{Wei2026}. Their construction supplies polynomial-resource procedures for logarithmic phase evolution. The present positive measure adds a separate Euclidean packet-depth evolution generated by $K_\beta$.

\begin{remark}
The positive Jacobi reduction can be interpreted as a Lanczos compression of the number-basis contraction. It produces a one-dimensional nearest-neighbor model that preserves a declared finite moment contract. This is a model-reduction statement. It is not a claim that the full number-theoretic quantum system is physically natural.
\end{remark}

\FlowChapter{The unresolved signed coherence}

For the principal zeta parameters define
\[
c_k
=
\sum_{n\ge1}
\frac{\mu(n)}{n^2}
\left(1-\frac1{n^2}\right)^k
\]
and
\[
E_k
=
\sum_{n\ge1}
\frac{\mu(n)^2}{n^4}
\left(1-\frac1{n^2}\right)^{2k}.
\]
The normalized coherence is
\[
\Gamma_k=\frac{\abs{c_k}^2}{E_k}.
\]
Since
\[
E_k\sim\frac{3}{2^{5/2}\pi^{3/2}}k^{-3/2},
\]
the classical B\'aez-Duarte criterion gives
\[
\boxed{
\mathrm{RH}
\iff
\Gamma_k=O_\varepsilon(k^\varepsilon)
\quad
\text{for every }\varepsilon>0.
}
\]
The positive results in this paper determine the denominator and its canonical reductions. They do not control the signed numerator.

\begin{boundary}[What remains]
Strict total positivity of $E_k$, positivity of the Jacobi model, the zero-free resolvent, and Laguerre universality do not imply subpolynomial M\"obius coherence. The missing estimate must use the arithmetic signs of $\mu(n)$ or an equivalent prime-sensitive structure.
\end{boundary}

\FigureWide{21_status_boundary}{The exact boundary of the paper.}{The positive side is complete within the stated setting. The signed principal coherence remains the open arithmetic problem.}

\InlinePart{Worked calculations and computational record}

\FlowChapter{Worked calculation of a four-node reduction}

Let $m=4$. Lanczos applied to $(T,v)$ produces
\[
J_4=
\begin{pmatrix}
a_0&b_1&0&0\\
b_1&a_1&b_2&0\\
0&b_2&a_2&b_3\\
0&0&b_3&a_3
\end{pmatrix}.
\]
The four-node rule is exact for degrees zero through seven:
\[
E_r=E_0e_1^TJ_4^re_1,
\qquad
0\le r\le7.
\]
Diagonalizing $J_4=U\Lambda U^T$ yields
\[
E_r=E_0\sum_{j=1}^{4}U_{1j}^2\Lambda_j^r.
\]
Thus an infinite squarefree sum is replaced by four positive packets without changing the first eight moments.

\begin{algorithm}[Constructing an $m$-packet reduced resolvent]
\begin{enumerate}[label=\arabic*.]
\item Form the diagonal nodes $x_n$ and seed amplitudes $\sqrt{w_n}$.
\item Run $m$ Lanczos steps using only diagonal multiplication by $x_n$.
\item Assemble the tridiagonal Jacobi matrix $J_m$.
\item Diagonalize $J_m$ to obtain positive nodes and weights.
\item Return
\[
\Ecal_m(z)=E_0\sum_{j=1}^{m}\frac{\lambda_j}{1-zx_j}.
\]
\item Certify the moments through degree $2m-1$.
\end{enumerate}
\end{algorithm}

\FlowChapter{Numerical tables and verification}

\input{generated_tables/checks_summary.tex}

The tables use a squarefree sieve through $1{,}500{,}000$. Low-order moment values are checked independently through the exact zeta-ratio formula at high precision. The asymptotic experiments are finite numerical certificates. Their purpose is to verify the formulas and display convergence, not to replace the proofs.

\appendix

\chapter{Formula registry}

\begin{longtable}{L{0.28\textwidth}L{0.64\textwidth}}
\toprule
Object & Formula\\
\midrule
\endfirsthead
\toprule
Object & Formula\\
\midrule
\endhead
Packet node & $x_n=(1-n^{-\beta})^2$\\
Packet weight & $w_n=\mu(n)^2\abs{\chi(n)}^2n^{-2a}$\\
Moment & $E_k=\sum_nw_nx_n^k$\\
Resolvent & $\Ecal(z)=\sum_nw_n/(1-zx_n)$\\
Gamma shape & $\alpha=(2a-1)/\beta$\\
Energy constant & $A=\delta_q\Gamma(\alpha)2^{-\alpha}/\beta$\\
Tilted law & $\dd\nu_k=x^k\dd\nu/E_k$\\
Scaled coordinate & $u=k(1-x)$\\
Laguerre polynomial & $\Lambda_m^{(\alpha)}=(-1)^mm!L_m^{(\alpha-1)}$\\
Laguerre norm & $h_m=m!(\alpha)_m$\\
Shifted determinant & $\Delta_d(k)=\det[E_{k+i+j}]_{i,j=0}^{d}$\\
Determinant constant & $G_d(\alpha)=\prod_{j=0}^{d}j!(\alpha)_j$\\
Principal cusp & $\Ecal(1)-\Ecal(z)\sim3(1-z)^{1/2}/(2^{3/2}\pi)$\\
RH coherence & $\Gamma_k=\abs{c_k}^2/E_k$\\
\bottomrule
\end{longtable}

\chapter{Implementation map}
The complete implementation is supplied as
\begin{center}
\texttt{projector\_generated\_special\_functions\_iii.py}.
\end{center}
It contains the squarefree sieve, packet-measure construction, high-precision zeta-ratio checks, diagonal Lanczos routine, orthogonal-polynomial reconstruction, determinant calculations, quadrature calculations, figure generation, table generation, and the built-in certificate ledger.

The central computational path is:
\begin{lstlisting}
params = PacketParameters(a=2.0, beta=2.0, conductor=1)
measure = build_measure(params, nmax=1_500_000)
energy = measure.energy(k)
global_model = global_lanczos(measure, steps=m)
tilted_model = tilted_lanczos(measure, k=k, steps=m)
nodes, weights = gauss_quadrature(tilted_model)
\end{lstlisting}

The Lanczos routine never materializes a dense operator. It multiplies vectors by the diagonal node array, orthogonalizes the resulting Krylov vectors, and returns the tridiagonal Jacobi coefficients. The source file is included unchanged in both release archives.

\chapter{Reproduction}
From the source directory, run
\begin{lstlisting}[language=bash,numbers=none]
python projector_generated_special_functions_iii.py
python generate_tables.py
python -m pytest -q
pdflatex -interaction=nonstopmode Projector_Generated_Special_Functions_III.tex
pdflatex -interaction=nonstopmode Projector_Generated_Special_Functions_III.tex
\end{lstlisting}

\begin{thebibliography}{20}
\bibitem{KomorebiSF1}
Komorebi P.,
\emph{Projector-Generated Special Functions I: Polyphase Differential Lifts, Bernoulli--Lerch Wavelets, Cyclotomic Detail Functions, and M\"obius Coherence}, 2026.

\bibitem{KomorebiSF2}
Komorebi P.,
\emph{Projector-Generated Special Functions II: Cyclotomic Stieltjes Packets, Packet Resultants, Positive Energy Resolvents, Gamma Scaling, and Character Decoherence}, 2026.

\bibitem{Akhiezer}
N. I. Akhiezer,
\emph{The Classical Moment Problem and Some Related Questions in Analysis},
Oliver and Boyd, 1965.

\bibitem{Chihara}
T. S. Chihara,
\emph{An Introduction to Orthogonal Polynomials},
Gordon and Breach, 1978.

\bibitem{Gautschi}
W. Gautschi,
\emph{Orthogonal Polynomials: Computation and Approximation},
Oxford University Press, 2004.

\bibitem{GolubWelsch}
G. H. Golub and J. H. Welsch,
Calculation of Gauss quadrature rules,
\emph{Mathematics of Computation} 23 (1969), 221--230.

\bibitem{Szego}
G. Szeg\H{o},
\emph{Orthogonal Polynomials},
American Mathematical Society, fourth edition, 1975.

\bibitem{DLMF}
NIST Digital Library of Mathematical Functions,
Chapters 18 and 25,
\url{https://dlmf.nist.gov/}.

\bibitem{BaezDuarte}
L. B\'aez-Duarte,
A sequential Riesz-like criterion for the Riemann hypothesis,
\emph{International Journal of Mathematics and Mathematical Sciences} 2005 (2005), 3527--3537.

\bibitem{Wei2026}
S. Wei, Y. Zhai, Q. Lu, et al.,
The Riemann Hypothesis manifested in dynamical quantum phase transitions,
\emph{Nature Communications}, accepted 18 June 2026,
DOI: 10.1038/s41467-026-74935-8.
\end{thebibliography}

\end{document}
